\begin{table}[H]
\centering
\caption{Nama Daerah Perbandingan Klaster STATIS Tahunan vs IPM (UNDP) (2020)} \label{tab:sample1statis_ipm_name_2020}
\begin{adjustbox}{max width=\textwidth}
\begin{tabular}{|>{\centering\arraybackslash}p{4cm}|p{10cm}|}
\hline
\multicolumn{1}{|c|}{\textbf{Kategori}} & \multicolumn{1}{c|}{\textbf{Daftar Wilayah}} \\
\hline
Cl. 1 (Rendah) \& Tinggi & Kota Tasikmalaya \\
\hline
Cl. 2 (Tinggi) \& Sangat Tinggi & Kota Jakarta Pusat \\
\hline
Cl. 2 (Tinggi) \& Tinggi & Kota Cilegon \\
\hline
Cl. 2 (Tinggi) \& Sedang & Bandung Barat \\
\hline
\end{tabular}
\end{adjustbox}
\end{table}